
\hypertarget{invited-alan-knee}{}
\label{invited-page-alan-knee}
\section{Plaskett Medal : Alan Knee}
\textbf{Session:}{ Tuesday -- Morning}\\
\textbf{Room:} A-1502\\

\paragraph{Biography}
{\sffamily
Dr. Alan Knee received his PhD in 2025 from the University of British Columbia, where he worked under the supervision of Dr. Jess McIver. He is currently a postdoctoral researcher at the University of Michigan. His thesis, "Searches for novel gravitational-wave sources with ground- and space-based detectors," was awarded the 2026 J.S. Plaskett Medal by CASCA for the most outstanding doctoral thesis in astronomy or astrophysics in Canada. His work enables novel insights into the ultra-dense matter of neutron stars, the dynamics of the Galactic centre, and the origin of distant LIGO-Virgo sources, requiring mastery of both ground- and space-based detector instrumentation and cutting-edge statistical approaches.
}

\paragraph{Abstract}
{\itshape
Searches for Novel Gravitational-Wave Sources with Ground- and Space-Based Detectors

The quest to expand the catalogue of gravitational-wave sources has driven the development of new data-analysis methods spanning a broad range of astrophysical targets and detector technologies. This thesis presents a set of independent but complementary searches for undiscovered gravitational-wave phenomena, from our own Galactic neighbourhood to the far reaches of the Universe. Topics explored include continuous gravitational waves from neutron stars in the Galactic centre, sub-threshold transient signals in LIGO-Virgo data, and the projected sensitivity of space-based detectors to exotic compact binaries. Together these works illustrate how pushing multiple frontiers simultaneously — source type, detector technology, and statistical methodology — can open new windows on the most energetic processes in the cosmos.
}


\hypertarget{invited-samantha-lawler-prized}{}
\label{invited-page-samantha-lawler-prized}
\section{Qilak Prize : Aaron Boyle and Samantha Lawler}
\textbf{Session:} {Wednesday -- Morning}\\
\textbf{Room:} A-1502\\
\textbf{Theme:} Night-Sky Protection / Solar System\\

\paragraph{Biography}
{\sffamily


Aaron Boley is a professor at the University of British Columbia and the co-director of the Outer Space Institute. His work on space sustainability is regularly featured in national and international news outlets, and he is a frequent guest on radio and TV, discussing astronomy and space exploration. Boley’s work on dark and quiet skies and space sustainability include engagement with delegations at the United Nations and policy makers at national and international agencies.

Samantha Lawler is an associate professor at the University of Regina. Her discoveries in the Kuiper Belt and predictions for satellite pollution have been featured by BBC, CBC, CNN, NPR, Scientific American, The New York Times, The Los Angeles Times, Wired Magazine, Nature, and many other international news outlets. She lives on a farm outside Regina and deeply appreciates the beautiful prairie skies.
}

\paragraph{Abstract}
{\itshape
As all Canadian astronomers and astronomy enthusiasts know, the night sky is under threat from light and radio-spectrum pollution, and from the overcrowding of orbits by satellites. In the face of these pressures, it has become crucial to cultivate understanding among both the public and lawmakers about these threats to the practice of astronomy. This is the important task that Profs. Lawler and Boley have undertaken. On issues related to light pollution, the proliferation of satellites and their effects on the night sky, radio spectrum, and orbital environments, they have met with government officials and decision-makers, and they have relentlessly petitioned the Canadian Space Agency, Innovation, Science, and Economic Development Canada, Transport Canada, the Department of National Defence, and Global Affairs Canada, as well as the U.S. Federal Communication Commission, the UK Space Agency and the UK Civil Aviation Authority. They have also made numerous appearances in the media, produced educational material for schools and presented dozens of public lectures to raise public awareness about these important issues.
CASCA is delighted to recognize Prof. Samantha Lawler and Prof. Aaron Boley for these outstanding contributions to the safeguarding of the night sky.
}
